% Section 3.1, from lectures/L03/appendix/a_interrupts.md.

\section{Interrupts and the Vector Table}
\label{c3:sec:interrupts}

\subsection{What an interrupt is on this machine}
\label{c3:sec:what}
An interrupt is a jump you did not write, taken between two instructions you did.

That is the whole of it. There is no scheduler, no thread, no context to switch. Some piece of
hardware sets a flag, the CPU finishes whatever instruction it was executing, pushes the program
counter, and jumps to a fixed address. Your handler runs, executes \code{reti}, and the program
counter comes back off the stack. The interrupted code has no way to know it happened.

Everything difficult about interrupts follows from that last sentence. Code that can be suspended
at any instruction boundary and resumed later is code you can no longer reason about one line at a
time, and \secref{c3:sec:contract} is about what to do instead.

\subsection{The vector table}
\label{c3:sec:table}
The first 52 words of program memory are not yours. They are a table of 26 slots, one per
interrupt source, and the hardware jumps to a fixed slot for each.

\bookfigure{vector_table}{The 26 ATmega328P interrupt vectors with their word addresses, from
RESET at \code{0x0000} to SPM\_READY at \code{0x0032}, with the ones this book uses
marked.}{c3:fig:vectors}

Four things about that table are worth having straight before you write anything.

\textbf{RESET is in it.} Vector 0 is the address the machine starts at, which is why a program's
first instruction has to be a jump to your code rather than the beginning of it: everything from
word \code{0x0002} upwards belongs to the table.

\textbf{Each slot is two words, so vector $n$ is at word $2n$.} PCINT0 is vector 3 and lives at
word \code{0x0006}, not \code{0x0003}. Two words, because a \code{jmp} on a part with 32\,KB of
flash is a two-word instruction, and the table has to hold one even though a hand-written table
usually puts a one-word \code{rjmp} there and leaves the second word unused.

\begin{aside}
\textbf{This book numbers the slots from zero}, so that RESET is vector 0 and the arithmetic is
$2n$ with nothing added. The datasheet's own table numbers them from \textbf{one}, so what it calls
``Vector No.~4'' is PCINT0, which is vector 3 here. The addresses are the same either way and only
the index differs; the offset is worth knowing before you check one against the other.
\end{aside}

\textbf{A slot holds a jump, not a handler.} Two words is room for one instruction and nothing
else, so what goes in a slot is \code{rjmp} to wherever your handler actually is. There is one
exception, and it is a useful trick: a handler that fits in a single instruction can go directly
in its slot, and \code{reti} alone is a legitimate handler that does nothing but acknowledge the
interrupt.

\textbf{The units already agree.} The datasheet's table is in words and AVRASM2 counts \code{.org}
in words, so the datasheet's number is the number you write: PCINT0 is at \code{0x0006} in the
table and \code{.org 0x0006} in your program. Nothing converts.

Better still, you need not write the number at all. \code{m328Pdef.inc} names every vector:

\begin{asm}
.org PCI0addr
    rjmp isr_pcint0
\end{asm}

\code{PCI0addr} is \code{0x0006}, \code{WDTaddr} is \code{0x000c}, \code{OC1Aaddr} is
\code{0x0016}, and the device file is the one place any of them is written down. A table built
from these names cannot drift from the datasheet, because it \emph{is} the datasheet.

\begin{aside}
Assembly written for \code{avr-gcc} counts \code{.org} in \textbf{bytes}, so the same PCINT0 slot
is reached by \code{.org 0x000C} there. If you are reading AVR assembly from elsewhere and every
handler seems to sit at twice its address, that is the dialect you are looking at, not a bug
(\secref{c6:sec:asmfromc} puts the two side by side).
\end{aside}

\textbf{And a C program does not write a table at all.} Where \code{avr-gcc} links the program,
avr-libc supplies the whole vector table and a macro fills one slot in:

\begin{cppcode}
ISR(TIMER1_COMPA_vect) { /* your handler */ }
\end{cppcode}

That builds the same two words at the same address. The names are the difference: avr-libc spells
the slots \code{TIMER1_COMPA_vect} and \code{PCINT0_vect} where \code{m328Pdef.inc} spells them
\code{OC1Aaddr} and \code{PCI0addr}, and one file never has both. \code{ISR} also writes the
handler's prologue and epilogue for you, saving what the compiler used and ending in \code{reti},
which is convenient right up to the point where you need to know exactly what was saved.
\code{ISR(vect, ISR_NAKED)} gives you the slot and nothing else: no prologue, no epilogue, no
\code{reti}, and every one of \secref{c3:sec:sreg}'s rules back in your hands. That is the form a
kernel's tick handler takes, and it is worth recognising now even though nothing in this book links
a C program to a vector.

\textbf{Build this table yourself before reading it again}, in both units. The word address is the
vector number times two, and the byte address is twice that again; the byte addresses are what the
simulator's program counter holds, so you will want them the moment a test reports where something
actually landed.

\subsection{The one flag that controls all of them}
\label{c3:sec:iflag}
Bit 7 of \code{SREG}, the \code{I} bit, is the global interrupt enable (\secref{c1:sec:sreg}). With
it clear, no interrupt is served, whatever else is configured. \code{sei} sets it and \code{cli}
clears it.

Enabling an interrupt therefore takes \textbf{two} switches, and forgetting either produces a
program that runs perfectly and never interrupts:
\begin{itemize}
  \item the peripheral's own enable, which for a pin change interrupt is two more bits still
    (\secref{c3:sec:pcint});
  \item the global \code{I} bit.
\end{itemize}

\code{I} is not a convenience. The hardware clears it on entry to any handler and \code{reti} sets
it again, which is what stops a handler from being interrupted by the same interrupt that started
it, and what makes the blocked window in \secref{c3:sec:blocked} a real quantity rather than a
worry.

\subsection{What happens, in order}
\label{c3:sec:sequence}

\bookfigure{interrupt_sequence}{The interrupt sequence as seven stacked steps, from the pin
changing, through finishing the current instruction, pushing the program counter, the vector jump,
the handler and \code{reti}, to one guaranteed main-program instruction, with the datasheet's
cycle cost beside each.}{c3:fig:sequence}

Two of those steps deserve more than the box gives them.

\textbf{``Finish the current instruction'' is why latency is a range rather than a number.} The
hardware does not abandon an instruction half-done, so if the flag is set just as a four-cycle
\code{ret} begins, three more cycles pass before anything else happens. The shortest possible
response is therefore \textbf{6 cycles} with an \code{rjmp} in the table, and the longest is
\textbf{9}.

\textbf{Both of those numbers assume an \code{rjmp} in the slot}, which is what this book writes.
Four cycles for the hardware's push, plus two for the jump, plus nought to three for the
instruction being finished. A table built with \code{jmp}, which is what \code{avr-gcc} emits
because it has to reach anywhere in flash, costs a cycle more at each end and the range is
\textbf{7 to 10}. The vector jump is a choice you make, so the entry latency is partly yours, and
quoting one range without saying which table it belongs to is how the two figures end up in the
same argument.

\textbf{``One more instruction, then the next'' is a guarantee, not an accident.} After
\code{reti}, the AVR always executes at least one instruction of the interrupted program before
serving another pending interrupt. Without that rule, a program under a steady stream of
interrupts could make no progress at all while appearing to run.

\subsubsection{What the simulator is and is not good for}
Every other cycle count in this book was measured. These were not, and the reason is worth
stating plainly rather than hiding.

simavr gets half of this right. It finishes the instruction it was executing before it takes an
interrupt, exactly as the device does: a flag that is set during a four-cycle \code{ret} waits up
to three cycles in the simulator as it would on the chip, which is also why Chapter~5's own
cross-check can explain a gap that alternates between 5332 and 5334 cycles by the two-cycle
\code{rjmp} the flag had to wait for. What it does not do is charge for the entry itself. It
pushes the program counter and jumps to the vector in no time at all, where the datasheet says the
push takes \textbf{4 cycles}, so every interrupt in the simulator reaches its handler four cycles
sooner than it would on the device. For that one quantity the datasheet is the authority and the
simulator is not, and the answer is arrived at by reading rather than by measuring.

That is not a complaint about simavr, which is exact about the thing this book leans on hardest:
what an instruction costs. It is a reminder that a tool is trustworthy about specific things, and
that ``I measured it'' is only an argument when you know the tool models what you measured.

\subsection{Pin change interrupts}
\label{c3:sec:pcint}
The ATmega328P has two kinds of external interrupt, and the difference matters.

\code{INT0} and \code{INT1} are two specific pins, each with its own vector, and each can be told
to fire on a rising edge, a falling edge, any change, or a low level. They are precise and there
are two of them.

\textbf{Pin change interrupts} cover 23 pins, which is very nearly every I/O pin on the part,
including the two that \code{INT0} and \code{INT1} already have, so a pin can have both kinds.
They are a much blunter instrument:
\begin{itemize}
  \item \textbf{One vector per port}, not per pin. All eight pins of port B share \code{PCINT0}.
  \item \textbf{Any change fires it.} There is no edge selection: pressed and released both
    interrupt.
  \item \textbf{It does not tell you which pin.} The handler is given nothing at all beyond the
    fact that something on that port changed.
\end{itemize}

\bookfigure{pcint_registers}{\code{PCICR} and \code{PCMSK0} as their bits: \code{PCIE0} is the
whole-port enable, \code{PCINT4} the single pin, and \code{PCINT7} and \code{PCINT6} are the
crystal pins.}{c3:fig:pcint}

Enabling one pin takes both registers:

\begin{booktable}{@{}p{5em}p{4.2em}L@{}}
  \thead{Register} & \thead{Address} & \thead{What it does} \\
  \midrule
  \code{PCICR} & \code{0x68} & One bit per port: \code{PCIE0} for B, \code{PCIE1} for C,
    \code{PCIE2} for D. \\
  \code{PCMSK0} & \code{0x6B} & One bit per pin of port B. \code{PCMSK1} and \code{PCMSK2} are C
    and D. \\
\end{booktable}

Both are at \code{0x60} and above, which is \textbf{extended I/O}, so \code{in} and \code{out}
cannot reach them and \code{lds} and \code{sts} are the only way in (\secref{c1:sec:iowindow}).
This is the first place in the book where that distinction stops being trivia and starts being an
assembler error.

\textbf{So a handler for two buttons on one port has to ask each of them whether it was the one.}
That is exactly what \code{btn_pressed} is for, and it is why the driver stores a structure per
button rather than a bit per port.

And there is one question the handler cannot answer at all: \textbf{which way the pin went}. The
interrupt fires on press and on release, and by the time the handler reads the pin the answer it
gets is the pin's state now, which for a fast enough press may already be the other one. The usual
fix is to compare against the previous state, kept in a variable, which is a variable shared
between a handler and everything else, which is \secref{c3:sec:contract}.
